\section{Discussion}
\label{sec:discussion}

\subsection{What needs explaining}
The empirical pattern is sharp.
On large contiguous center holes, increasing RePaint's resampling count $r$ \emph{hurts} our mask-conditioned DDPM and \emph{helps} our unconditional DDPM (Table~\ref{tab:resampling-delta}, Figures~\ref{fig:cond-r} and~\ref{fig:uncond-r}).
Matched-coverage brush masks show no significant $r$ effect (Table~\ref{tab:brush-null}).
RePaint's own pretrained unconditional CelebA-HQ checkpoint likewise improves with $r$ (Figure~\ref{fig:external-repaint}).
We interpret this as a mismatch between RePaint's inference-time correction loop and how much measurement information the network already receives during each denoising step.

\subsection{Unconditional: resampling as alternating projections}
\label{sec:alt-proj}
In the classic RePaint setting the U-Net never sees a mask: $\varepsilon_\theta=\varepsilon_\theta(x_t,t)$.
Its reverse step proposes a full-image update $\hat{x}_{t-1}^{\mathrm{unk}}$; consistency with known pixels is enforced only by the stitch~\eqref{eq:stitch}, and when $r{>}1$ by forward jumps~\eqref{eq:jump}.
Writing $P_{\mathrm{prior}}(x_t)$ for one reverse (denoising) step that uses $\varepsilon_\theta(x_t,t)$, and $P_{\mathrm{data}}$ for the known-pixel stitch, a resampling cycle is the composition
\begin{align}
  x_{t-1}
  &\leftarrow
  \underbrace{m\odot q(x_0,t{-}1)+(1-m)\odot P_{\mathrm{prior}}(x_t)}_{\text{known-pixel projection }P_{\mathrm{data}}},
  \label{eq:proj-data}\\
  x_{t}
  &\leftarrow
  \underbrace{\sqrt{1-\beta_{t}}\,x_{t-1}+\sqrt{\beta_{t}}\,\varepsilon}_{\text{re-noise / jump}},
  \qquad \varepsilon\sim\mathcal{N}(0,I),
  \label{eq:proj-jump}
\end{align}
repeated within the RePaint $(j,r)$ schedule.
This is an instance of \emph{alternating projections}: $P_{\mathrm{prior}}$ moves toward the learned image manifold, $P_{\mathrm{data}}$ projects back onto the known-pixel constraint, and~\eqref{eq:proj-jump} reopens the hole so the pair can be applied again.
Without those projections, the hole and the context disagree at the boundary---hence the blotchy seams at $r{=}1$ and the clear gain from $r{=}5$ on our unconditional model (and on the external CelebA-HQ checkpoint).
Resampling is doing real work that the architecture does not otherwise provide.

\subsection{Conditioned: a static conditioning channel already carries the constraint}
Our conditioned model instead predicts $\varepsilon_\theta(x_t,t,\tilde{x}_0,m)$ from the concatenation $(x_t,\tilde{x}_0,m)$ at every reverse step.
Known content is therefore not only an after-the-fact constraint $P_{\mathrm{data}}$; it is a \emph{static conditioning channel} the network was trained to read while predicting $\varepsilon$.
Under that training regime, a single application of~\eqref{eq:stitch} at $r{=}1$ is usually enough: the reverse process is already steered by $(\tilde{x}_0,m)$, and large-center fills at $r{=}1$ are coherent (Figure~\ref{fig:cond-r}a).

Raising $r$ re-imposes the alternating-projection loop~\eqref{eq:proj-data}--\eqref{eq:proj-jump} on top of that channel.
Each extra jump re-noises the hole, re-denoises, and re-applies $P_{\mathrm{data}}$---repeatedly overwriting a fill that the conditioned predictor had already committed to.
We hypothesize that this second consistency mechanism is largely redundant for a mask-conditioned $\varepsilon_\theta$, and on large contiguous holes it is actively harmful: the hole is big enough that re-harmonization can replace a sharp, identity-specific completion with an over-smoothed / washed-out compromise (Figure~\ref{fig:cond-r}b).
In short, alternating projections fix a missing constraint in the unconditional case; they fight a constraint that is already inside the network in the conditioned case.

\subsection{Why the brush null fits}
If the harm were merely ``more compute'' or ``any missing pixels,'' matched-coverage brush masks should degrade similarly under $r{=}10$.
They do not (Table~\ref{tab:brush-null}).
Brush strokes leave dense local context adjacent to every missing pixel, so both the static conditioning channels and a single stitch already determine the fill; there is little contiguous region for repeated re-harmonization to rewrite.
The conditioned $r$-harm is therefore geometry-specific: it tracks large contiguous holes where RePaint's loop has a coherent interior to overwrite, not masked area alone.

\subsection{Practical implication}
RePaint's default $(j,r)$ was designed for unconditional generators.
Copying that recipe onto a Palette-style conditioned DDPM is not a neutral choice: on hard center masks it can cost ${\approx}10{\times}$ NFEs while lowering PSNR by ${\sim}2.7$\,dB.
A regime-aware rule is simpler---use moderate resampling for our unconditional model (and classic unconditional RePaint); prefer $r{=}1$ (noise-matched stitch only) for mask-conditioned models on large contiguous holes (brush/sparse masks are insensitive to $r$ either way, so $r{=}1$ is a safe default there too)---and report $(j,r)$ jointly with the training mode.

\subsection{Limitations of the mechanism story}
This account is an interpretation of paired ablations, not a proof of the internal computation.
We do not intervene on intermediate latents, ablate the conditioning channels at inference, or derive a formal fixed-point for the RePaint operator under concat-conditioning.
Alternative explanations remain open, including capacity limits that make repeated resampling amplify a bad mode.

\paragraph{Schedule granularity.}
A particularly natural alternative---raised by RePaint's own discussion of per-step harmonization budget under finer schedules---is that resampling benefits are schedule-specific rather than conditioning-specific.
Our headline grid (Table~\ref{tab:resampling-delta}) runs both of our models at full $T{=}1000$; the external check (Figure~\ref{fig:external-repaint}) uses a coarse respaced $T{=}100$ schedule but only for RePaint's pretrained CelebA-HQ checkpoint.
We therefore evaluate the conditioned model only at full $T{=}1000$; whether the resampling-harm effect weakens or disappears at coarser, respaced schedules (as in RePaint's original $T{=}250$ setting) remains untested and is a natural target for follow-up work.
That said, granularity alone does not explain the \emph{within}-schedule sign flip: at matched $T{=}1000$, the same $(j,r)$ recipe helps our unconditional model and hurts our conditioned one, which still points to conditioning as a driver even if schedule interacts with effect size.

What the present data do pin down is that sign flip, its concentration on contiguous holes, and the fact that RePaint's pretrained CelebA-HQ checkpoint does not show conditioned-style degradation under higher $r$.

\subsection{Scope limitations}
Beyond the mechanism caveats above, the study itself is narrow.
All quantitative claims use a single checkpoint epoch ($30$), one U-Net capacity recipe, and $64{\times}64$ CelebA only.
Paired tests use $n{=}32$ seed-matched images; effect sizes are large for the conditioned harm, but broader populations and datasets are untested.
Mask coverage is limited to center squares plus one brush control---not free-form or semantic masks at scale---and we do not ablate learned variance, classifier-free guidance, or latent-space / Stable-Diffusion-style inpainting pipelines.
Wall-clock numbers are hardware-specific; NFE ratios are the portable cost measure.
These bounds do not overturn the paired interaction we report, but they limit how far the practical ``$r{=}1$ for conditioned models, moderate $r$ for our unconditional model'' rule should be extrapolated without further evidence.
